% =====================================================================
\chapter{Funciones, imagen y preimagen}
\label{tema:2}
% =====================================================================

Este tema define con precisión qué es una función y estudia cómo actúa sobre subconjuntos mediante la imagen y la preimagen. Sigue las secciones~12.1 y~12.6 de Hammack~\cite{hammack}, las secciones~5.1 y~5.5 de Velleman~\cite{velleman} y la sección~4.3 de~\cite{mcs}. Las propiedades de la imagen y la preimagen son las de la pregunta~9 de la hoja~1 de Cambridge~\cite{camns}.

\section{Definición de función}

\begin{definicion}[Función]\label{def:t2-funcion}
Sean $A$ y $B$ conjuntos. Una \emph{función} $f$ de $A$ en $B$, escrita $f\colon A\to B$, es un subconjunto $f\subseteq A\times B$ tal que para cada $a\in A$ existe \emph{exactamente un} $b\in B$ con $(a,b)\in f$; ese $b$ se denota $f(a)$. El conjunto $A$ es el \emph{dominio} y $B$ el \emph{codominio}. El \emph{recorrido} (o \emph{imagen}) de $f$ es $f(A)=\{f(a): a\in A\}\subseteq B$.
\end{definicion}

La definición contiene dos exigencias: \emph{existencia} (todo elemento del dominio tiene imagen) y \emph{unicidad} (no tiene dos). Dos funciones $f,g\colon A\to B$ son iguales si $f(a)=g(a)$ para todo $a\in A$; la igualdad exige además que el dominio y el codominio coincidan.

\begin{ejemplo}\label{ej:t2-funciones}
\begin{enumerate}
  \item $f\colon\Z\to\Z$, $f(n)=n^{2}$, es una función con recorrido $\{0,1,4,9,\dots\}\subsetneq\Z$.
  \item La \emph{identidad} $\id_A\colon A\to A$, $\id_A(a)=a$.
  \item La \emph{longitud} $\ell\colon\Sigma^{*}\to\N$, $\ell(w)=\abs{w}$, asigna a cada cadena su número de símbolos.
  \item Las \emph{proyecciones} $p_1\colon A\times B\to A$, $p_1(a,b)=a$, y $p_2(a,b)=b$.
  \item La regla «$x\mapsto\sqrt{x}$» \emph{no} define una función $\R\to\R$: falla la existencia, porque $\sqrt{-1}\notin\R$. La regla «$x\mapsto y$ tal que $y^{2}=x$» de $[0,\infty)$ en $\R$ falla la unicidad, porque $4$ tendría las imágenes $2$ y $-2$.
\end{enumerate}
\end{ejemplo}

\begin{ejemplo}[Funciones bien definidas]\label{ej:t2-biendef}
Cuando los elementos del dominio admiten varias representaciones, hay que comprobar que la regla no depende de la representación elegida. La regla $f\bigl(\tfrac{a}{b}\bigr)=a+b$ no define una función $\Q\to\Z$: $\tfrac12=\tfrac24$, pero $1+2=3\neq 6=2+4$. En cambio, $g\bigl(\tfrac ab\bigr)=\tfrac{2a}{b}$ sí está bien definida: si $\tfrac ab=\tfrac cd$, entonces $\tfrac{2a}{b}=2\cdot\tfrac ab=2\cdot\tfrac cd=\tfrac{2c}{d}$.
\end{ejemplo}

\section{El conjunto de funciones}

\begin{definicion}
$B^{A}$ denota el conjunto de todas las funciones $A\to B$.
\end{definicion}

\begin{proposicion}\label{prop:t2-contar}
Si $A$ y $B$ son finitos, $\abs{B^{A}}=\abs{B}^{\abs{A}}$.
\end{proposicion}

\begin{proof}
Sea $A=\{a_1,\dots,a_m\}$. Una función $f\colon A\to B$ queda determinada por la lista de valores $(f(a_1),\dots,f(a_m))\in B^{m}$, y toda lista de $B^{m}$ define una función. Luego $\abs{B^{A}}=\abs{B^{m}}=\abs{B}^{m}$ por la \cref{prop:t1-producto}.
\end{proof}

\begin{definicion}[Función característica]
Sea $X\subseteq U$. La \emph{función característica} (o \emph{indicadora}) de $X$ es $\chi_X\colon U\to\{0,1\}$, con $\chi_X(u)=1$ si $u\in X$ y $\chi_X(u)=0$ si $u\notin X$.
\end{definicion}

Cada subconjunto de $U$ determina su función característica y cada función $U\to\{0,1\}$ es la característica de un único subconjunto, el de los elementos que van al~1. Con $\abs{U}=n$, la \cref{prop:t2-contar} da $\abs{\{0,1\}^{U}}=2^{n}$, de nuevo el \cref{teo:t1-partes}. En el \cref{tema:3} se formaliza esta correspondencia como una biyección.

\section{Imagen y preimagen de subconjuntos}

\begin{definicion}[Imagen y preimagen]\label{def:t2-imagen}
Sea $f\colon A\to B$. Para $X\subseteq A$ y $Y\subseteq B$ se definen la \emph{imagen} de $X$ y la \emph{preimagen} (o \emph{imagen inversa}) de $Y$:
\[
f(X)=\{f(x): x\in X\}\subseteq B,\qquad f^{-1}(Y)=\{a\in A: f(a)\in Y\}\subseteq A.
\]
\end{definicion}

La notación $f^{-1}(Y)$ tiene sentido para \emph{cualquier} función, sea o no invertible: es un conjunto, no la aplicación de una función inversa. Si $Y=\{b\}$, el conjunto $f^{-1}(\{b\})$ se llama \emph{fibra} de $b$.

\begin{ejemplo}\label{ej:t2-cuadrado}
Sea $f\colon\R\to\R$, $f(x)=x^{2}$. Entonces $f([-1,2])=[0,4]$, porque $x^{2}$ recorre $[0,1]$ en $[-1,0]$ y $[0,4]$ en $[0,2]$. Además, $f^{-1}([1,4])=[-2,-1]\cup[1,2]$, $f^{-1}(\{-1\})=\emptyset$ y $f^{-1}(f([0,1]))=f^{-1}([0,1])=[-1,1]\supsetneq[0,1]$.
\end{ejemplo}

\begin{teorema}[Propiedades de la preimagen]\label{teo:t2-preimagen}
Sea $f\colon A\to B$ y sean $C,D\subseteq B$. Entonces
\begin{gather*}
f^{-1}(C\cup D)=f^{-1}(C)\cup f^{-1}(D),\qquad f^{-1}(C\cap D)=f^{-1}(C)\cap f^{-1}(D),\\
f^{-1}(B\setminus C)=A\setminus f^{-1}(C).
\end{gather*}
\end{teorema}

\begin{proof}
Se demuestra la primera; las otras dos se dejan como ejercicio (\cref{eje:2-preim-ui,eje:2-preim-comp}). Para todo $a\in A$:
$a\in f^{-1}(C\cup D)\iff f(a)\in C\cup D\iff f(a)\in C\ \text{o}\ f(a)\in D\iff a\in f^{-1}(C)\ \text{o}\ a\in f^{-1}(D)\iff a\in f^{-1}(C)\cup f^{-1}(D)$.
\end{proof}

\begin{teorema}[Propiedades de la imagen]\label{teo:t2-imagen}
Sea $f\colon A\to B$, $X_1,X_2,X\subseteq A$ e $Y\subseteq B$. Entonces
\begin{enumerate}
  \item $f(X_1\cup X_2)=f(X_1)\cup f(X_2)$;
  \item $f(X_1\cap X_2)\subseteq f(X_1)\cap f(X_2)$, y la inclusión puede ser estricta;
  \item $X\subseteq f^{-1}(f(X))$ y $f(f^{-1}(Y))\subseteq Y$.
\end{enumerate}
\end{teorema}

\begin{proof}
(ii) Si $b\in f(X_1\cap X_2)$, existe $x\in X_1\cap X_2$ con $f(x)=b$; como $x\in X_1$, $b\in f(X_1)$, y como $x\in X_2$, $b\in f(X_2)$. Para la desigualdad estricta, sea $f(x)=x^{2}$, $X_1=\{-1\}$ y $X_2=\{1\}$: $f(X_1\cap X_2)=f(\emptyset)=\emptyset$, mientras que $f(X_1)\cap f(X_2)=\{1\}$. Los apartados (i) y~(iii) son los \cref{eje:2-imagen-ui,eje:2-ida-vuelta}.
\end{proof}

La asimetría entre ambos teoremas es instructiva: la preimagen respeta todas las operaciones de conjuntos, mientras que la imagen solo respeta la unión. La razón es que la imagen puede «pegar» elementos distintos; en el \cref{tema:3} se verá que las igualdades $f(X_1\cap X_2)=f(X_1)\cap f(X_2)$ y $f^{-1}(f(X))=X$ son ciertas exactamente cuando $f$ es inyectiva.

\section{Otras construcciones}

Una \emph{sucesión} de elementos de $A$ es una función $\N\to A$, $n\mapsto a_n$. Una \emph{cadena} de longitud $n$ sobre $\Sigma$ puede verse como una función $\{1,\dots,n\}\to\Sigma$ que asigna a cada posición su símbolo. Si $f\colon A\to B$ y $X\subseteq A$, la \emph{restricción} $f|_X\colon X\to B$ es la función $f|_X(x)=f(x)$. Las funciones \emph{definidas a trozos} asignan reglas distintas a partes disjuntas del dominio que lo recubren.

\begin{observacion}[Conexión con el curso]
Un modelo de lenguaje es una función $P\colon\V^{*}\to[0,1]$ (módulo~1, §1.3), y una matriz de \emph{embeddings} define una función $E\colon\V\to\R^{d}$ (módulo~12, tema~4). La preimagen aparece de forma natural: para la longitud $\ell\colon\V^{*}\to\N$, la fibra $\ell^{-1}(\{T\})$ es $\V^{T}$; para la función que elige el token de mayor logit (decodificación voraz, módulo~1, tema~4), la preimagen de un token es la región del espacio de logits en la que ese token gana.
\end{observacion}

\begin{observacion}[Errores frecuentes]
(1) Confundir codominio y recorrido: $f(x)=x^{2}$ de $\R$ en $\R$ tiene codominio $\R$ y recorrido $[0,\infty)$. (2) Leer $f^{-1}(Y)$ como «aplicar la inversa»: la preimagen existe siempre, la inversa solo para biyecciones. (3) Olvidar comprobar que una regla está bien definida cuando el dominio tiene representaciones múltiples (fracciones, clases de equivalencia). (4) Suponer que $f^{-1}(f(X))=X$: solo se tiene la inclusión $\supseteq$ en general.
\end{observacion}

\begin{lecturas}
\item Hammack~\cite{hammack}, §12.1 «Functions» y §12.6 «Image and Preimage».
\item Velleman~\cite{velleman}, §5.1 «Functions» y §5.5 «Images and Inverse Images: A Research Project».
\item Lehman, Leighton y Meyer~\cite{mcs}, §4.3 «Functions».
\end{lecturas}

\section*{Ejercicios del tema 2}
\addcontentsline{toc}{section}{Ejercicios del tema 2}

\begin{ejercicio}{1}{eje:2-esfuncion}
Decida si cada una de las siguientes es una función y, si no lo es, qué exigencia falla. (a) $R_1=\{(1,\cad{a}),(2,\cad{a}),(3,\cad{b})\}\subseteq\{1,2,3\}\times\{\cad{a},\cad{b}\}$. (b) $R_2=\{(1,\cad{a}),(2,\cad{b})\}$, con los mismos conjuntos. (c) $R_3=\{(1,\cad{a}),(1,\cad{b}),(2,\cad{a}),(3,\cad{a})\}$. (d) $x\mapsto 1/x$ de $\R$ en $\R$. (e) $\tfrac ab\mapsto\tfrac{a+1}{b+1}$ de $\Q_{>0}$ en $\Q$, con $a,b\in\N$, $a,b\ge1$.
\end{ejercicio}

\begin{ejercicio}{2}{eje:2-tabla}
Sean $A=\{0,1,2,3,4\}$, $B=\{1,2,3,4,5\}$ y
\[
f=\{(0,2),(1,4),(2,2),(3,1),(4,2)\}.
\]
(a) Compruebe que $f$ es una función $A\to B$ y dé su dominio, codominio y recorrido. (b) Calcule $f(\{0,1,2\})$, $f^{-1}(\{2\})$, $f^{-1}(\{3,5\})$ y $f^{-1}(\{1,4\})$.
\fuente{adaptado de Hammack~\cite{hammack}, §12.1 «Functions», ejercicios de la sección (datos cambiados).}
\end{ejercicio}

\begin{ejercicio}{2}{eje:2-contar}
(a) ¿Cuántas funciones hay de un conjunto de 3 elementos en uno de 4? ¿Y de uno de 4 en uno de 3? (b) ¿Cuántas funciones $\{0,1\}^{2}\to\{0,1\}$ (funciones booleanas de dos variables) existen? (c) ¿Cuántas funciones $\{0,1\}^{n}\to\{0,1\}$? Evalúe para $n=3$.
\end{ejercicio}

\begin{ejercicio}{2}{eje:2-cuadrado}
Sea $f\colon\R\to\R$, $f(x)=x^{2}-2x$. Calcule $f([0,3])$, $f^{-1}(\{3\})$, $f^{-1}([-1,0])$, $f^{-1}((-\infty,-2))$ y $f^{-1}(f(\{3\}))$.
\fuente{adaptado de Hammack~\cite{hammack}, §12.6 «Image and Preimage», ejercicios de la sección.}
\end{ejercicio}

\begin{ejercicio}{2}{eje:2-suelo}
Sea $f\colon\Z\to\Z$, $f(n)=\lfloor n/2\rfloor$. Calcule $f(\{0,1,2,3\})$, $f^{-1}(\{1\})$, $f^{-1}(\{0,1,2\})$ y $f(2\Z)$, donde $2\Z$ es el conjunto de los enteros pares. ¿Cuál es el recorrido de $f$?
\end{ejercicio}

\begin{ejercicio}{2}{eje:2-longitud}
Sea $\Sigma=\{\cad{a},\cad{b}\}$ y $\ell\colon\Sigma^{*}\to\N$ la longitud. Calcule $\ell^{-1}(\{2\})$, $\ell^{-1}(\{0,1\})$, $\ell(\{\cad{ab},\cad{ba},\cad{aab}\})$ y el recorrido de $\ell$. ¿Cuántos elementos tiene $\ell^{-1}(\{n\})$?
\end{ejercicio}

\begin{ejercicio}{1}{eje:2-recorrido}
Dé el recorrido de cada función $\R\to\R$: (a) $f(x)=e^{x}$; (b) $g(x)=\operatorname{sen}x$; (c) $h(x)=x^{2}+2x+3$; (d) $k(x)=\dfrac{1}{1+x^{2}}$.
\end{ejercicio}

\begin{ejercicio}{3}{eje:2-preim-ui}
Sea $f\colon X\to Y$ y $C,D\subseteq Y$. Demuestre que $f^{-1}(C\cap D)=f^{-1}(C)\cap f^{-1}(D)$ y que, si $C\subseteq D$, entonces $f^{-1}(C)\subseteq f^{-1}(D)$.
\fuente{adaptado de Cambridge, IA Numbers and Sets~\cite{camns}, Example Sheet~1, pregunta~9.}
\end{ejercicio}

\begin{ejercicio}{3}{eje:2-imagen-ui}
Sea $f\colon X\to Y$ y $A,B\subseteq X$. Demuestre que $f(A\cup B)=f(A)\cup f(B)$. Decida si $f(A\setminus B)=f(A)\setminus f(B)$ es siempre cierta; si no, dé un contraejemplo y la inclusión que sí se cumple.
\fuente{adaptado de Cambridge, IA Numbers and Sets~\cite{camns}, Example Sheet~1, pregunta~9.}
\end{ejercicio}

\begin{ejercicio}{3}{eje:2-ida-vuelta}
Sea $f\colon X\to Y$, $A\subseteq X$ y $C\subseteq Y$. Demuestre que $A\subseteq f^{-1}(f(A))$ y que $f(f^{-1}(C))\subseteq C$. Dé ejemplos en los que ambas inclusiones son estrictas.
\fuente{adaptado de Cambridge, IA Numbers and Sets~\cite{camns}, Example Sheet~1, pregunta~9.}
\end{ejercicio}

\begin{ejercicio}{3}{eje:2-preim-comp}
Sea $f\colon X\to Y$ y $C\subseteq Y$. Demuestre que $f^{-1}(Y\setminus C)=X\setminus f^{-1}(C)$, que $f^{-1}(\emptyset)=\emptyset$ y que $f^{-1}(Y)=X$.
\fuente{adaptado de Cambridge, Discrete Mathematics~\cite{camdm}, hoja~8, ejercicio~12.2.5.}
\end{ejercicio}

\begin{ejercicio}{2}{eje:2-caracteristica}
Sean $X,Z\subseteq U$. Compruebe, distinguiendo casos según la pertenencia de $u$ a $X$ y a $Z$, que para todo $u\in U$:
$\chi_{X\cap Z}(u)=\chi_X(u)\chi_Z(u)$, \ $\chi_{X\cup Z}(u)=\chi_X(u)+\chi_Z(u)-\chi_X(u)\chi_Z(u)$ \ y \ $\chi_{\overline{X}}(u)=1-\chi_X(u)$.
Deduzca que, si $U$ es finito, $\abs{X\cup Z}=\abs{X}+\abs{Z}-\abs{X\cap Z}$.
\end{ejercicio}

\begin{ejercicio}{3}{eje:2-biendef}
Decida si las reglas siguientes definen funciones $\Q\to\Q$ y demuéstrelo: (a) $f\bigl(\tfrac ab\bigr)=\tfrac{a^{2}}{b^{2}}$; (b) $g\bigl(\tfrac ab\bigr)=\tfrac{a+b}{b}$; (c) $h\bigl(\tfrac ab\bigr)=\tfrac{a}{b+1}$ (con $b\ge1$). En todos los casos $a\in\Z$ y $b\in\Z$, $b\neq0$.
\end{ejercicio}

\begin{ejercicio}{2}{eje:2-argmax}
Sea $g\colon\R^{3}\to\{1,2,3\}$ la función que devuelve el índice de la mayor coordenada y, en caso de empate, el menor de los índices empatados (decodificación voraz). (a) Calcule $g(1,3,2)$, $g(2,2,1)$ y $g(0,0,0)$. (b) Describa $g^{-1}(\{1\})$ y $g^{-1}(\{3\})$ mediante desigualdades. (c) ¿Es $g^{-1}(\{1\})\cup g^{-1}(\{2\})\cup g^{-1}(\{3\})=\R^{3}$? ¿Son disjuntas las tres fibras?
\end{ejercicio}
